x\section{Polynomials}
\subsection{Introduction to Polynomials}
\begin{definition}[Polynomial Function]
  A polynomial function, in general, is also stated as a polynomial or polynomial expression, defined by its degree. The degree of any polynomial is the highest power present in it.

  $$f(x) = a_nx^n + a_{n-1}x^{n-1} + ... + a_2x^2 + a_1x + a_0$$
  $$\text{Leading Term: } a_nx^n$$
  $$\text{Constant Term: } a_0$$
  $$\text{Degree: }n \text{ \textit{(Largest exponent in the polynomial.)}}$$


  The domain of a polynomial is all $\R$ numbers.
\end{definition}

\textbf{Note:} Exponents are all non negative integers.

\begin{example}[Polynomial Function]
  \hfill
  \begin{enumerate}
  \item $f(x) = 5x^2 + 4x^4$ \\
    $f(x) = 4x^4 + 5x^2$ \\
    This is a degree 4 polynomial.
  \item $g(x) = 3-\frac{1}{2}x$ \\
    $-\frac{1}{2}x + 3$ \\
    This is a degree 1 polynomial.
  \item $h(x) = 9$ \\
    This is a degree 0 polynomial.
  \item $f(x) = 0$ \\
    This function doesn't have a degree.
  \item $f(x) = 9x^4 - \pi x^2 + \frac{1}{2}$
  \item $h(x) = \frac{3x^2 - 1}{7}$ \\
    $\frac{3x^2}{7} - \frac{1}{7}$ \\
    $\frac{3}{7}x^2 - \frac{1}{7}$
  \item $F(x) = 2x^3(x^2 - 1)(x - 1)^2$ \\
    \textbf{Note: }We can do a fake distribution here. \\
    $2 \cdot x^3 \cdot (x^2 - 1)^1 \cdot (x^1 - 1)^2$ \\
    $2 \cdot x^3 \cdot x^2 \cdot x^2$ \\
    $2x^7$ 
  \end{enumerate}
\end{example}

\begin{example}[Non-Polynomial Function]
  \hfill
  \begin{enumerate}
  \item $F(x) = x^\frac{2}{3} + 2x - 1$
  \item $G(x) = 1 - \frac{4}{x}$
  \item $g(x) = \frac{x^2 - 5}{x^4}$ \\
    $\frac{x^2}{x^4} - \frac{5}{x^4}$ \\
    $\frac{1}{x^2} - \frac{5}{x^4}$
  \item $G(x) = 2\sqrt{x}(\sqrt{x} + 1)$
  \end{enumerate}
\end{example}
